\makeatletter
\def\input@path{{./}{../../../preamble/}{book2026-2027/preamble/}}
\makeatother
\input{preamble.tex}
\usepackage{needspace}
\hypersetup{pdftitle={Grade 10 Term 2 Discrete Probability}}
\setlength{\emergencystretch}{2em}
\title{Grade 10, Term 2, Lesson 2, Criterion C7\\Poisson Event Counts and Exposure}
\author{}
\date{}
\begin{document}
\maketitle
\setcounter{tocdepth}{1}
\setcounter{secnumdepth}{2}
\phantomsection\label{toc}
\tableofcontents
\section*{Evaluated criterion}
\textbf{C7:} Evaluates rare-event situations with the Poisson distribution, interpreting the parameter $\lambda$ in contexts such as claims, arrivals, defects, or defaults within a fixed interval.
\section{Introduction}\label{sec:t2-c7-introduction}
All data in this activity are hypothetical teaching models. Probabilities have no monetary units. Keep exact expressions where useful and round final probabilities to four decimal places.
\subsection{Count events over an exposure}
A school shop receives occasional refund requests. Observe requests during one hour rather than fixing customer count. Let $X$ count requests in that hour. A Poisson variable has support $\{0,1,2,\ldots\}$ events: zero is possible and there is no fixed largest count.
\textbf{Check.} Does this variable measure time until the next request? \emph{Suggested answer:} No. It counts events over a stated interval.


\subsection{Rate and interval mean}
A rate of 0.5 requests per hour gives expected count $\lambda=0.5$ for one hour and $\lambda=1$ for two hours:
\[\lambda=\text{rate}\times\text{exposure}.\]
The Greek letter $\lambda$, read lambda, is the expected count in the chosen exposure. A rate in events per hour times hours gives events. The expected count may be non-integer, although observed counts are integers.
\textbf{Check.} At 2 claims per day, what mean applies to half a day? \emph{Suggested answer:} $\lambda=2(0.5)=1$ claim.


\subsection{Assumptions and notation}
Assume constant event rate, independent counts in disjoint intervals, and isolated events: in a very short interval the chance of one event is approximately proportional to its length and that of two or more is negligible. These assumptions do not follow from rarity. A shared storm can cluster claims. The constant $e\approx2.71828$ is a mathematical number on a calculator. Factorial means $k!=k(k-1)\cdots1$ and $0!=1$; for example $3!=6$.
\textbf{Check.} Why is $k=0$ allowed? \emph{Suggested answer:} An interval can contain no events; $0!=1$.

\textbf{Check.} Does rarity alone establish Poisson? \emph{Suggested answer:} No. Check rate stability, independent intervals and absence of clustering.

\section{General worked problems}
\Needspace{14\baselineskip}
\subsection{Introductory general problem: two events in one unit}
An ideal process has constant rate 1 event per exposure unit, independent disjoint-interval counts and isolated events.

\textbf{Questions/Tasks}
\begin{enumerate}
\item Calculate $P(X=2)$ for the exposure. 
\end{enumerate}
\subsubsection*{Complete worked solution}
\begin{center}
\begin{tabular}{llp{0.48\linewidth}}
\toprule
Symbol & Value & Meaning in this problem \\ \midrule
rate & 1 event per unit & Event frequency per exposure unit \\
exposure & 1 unit & Observation length \\
$\lambda$ & 1 event & Expected count for this exposure \\
\bottomrule
\end{tabular}
\end{center}

\par\noindent\begin{minipage}{\linewidth}
\textbf{Step 1: Define and identify units.} $X$ counts events over 1 exposure unit; support $\{0,1,2,\ldots\}$ events. The count differs from rate and money.
\end{minipage}\par\medskip
\par\noindent\begin{minipage}{\linewidth}
\textbf{Step 2: Check and scale.} Under constant rate, independent disjoint-interval counts and isolated events, use Poisson. $\lambda=\text{rate}\times\text{exposure}=1\times1=1$ events.
\end{minipage}\par\medskip
\par\noindent\begin{minipage}{\linewidth}
\textbf{Step 3: Translate and calculate.} The event is $ X=2$. Symbolically,
\[P(X=k)=e^{-\lambda}\lambda^k/k!.\]
Numerically,
\[P(X=2)=e^{-1}(1)^2/2!\approx 0.1839.\]
\end{minipage}\par\medskip
\par\noindent\begin{minipage}{\linewidth}
\textbf{Step 4: Check and interpret.} The probability is in $[0,1]$; the full Poisson PMF sums to 1. 
\end{minipage}\par\medskip
\subsection{Bridge to the general model}
\[X\sim\operatorname{Pois}(\lambda),\qquad P(X=k)=\frac{e^{-\lambda}\lambda^k}{k!},\quad k=0,1,\ldots.\]
The numerator and factorial assign the mass at $k$. Use the mean for the interval in the question, not automatically the reported rate. If $\lambda=0$, no events occur with probability 1; for $\lambda>0$, support is all non-negative integers. The sum $\sum_{j=0}^{b}$ adds counts 0 through $b$; at least $a$ complements at most $a-1$.
\Needspace{14\baselineskip}
\subsection{Advanced general problem: rescaling and a complement}
An ideal process has constant rate 0.8 events per hour, independent disjoint-interval counts and isolated events. Observe 2.5 hours.

\textbf{Questions/Tasks}
\begin{enumerate}
\item Calculate $P(X\geq3)$ for the exposure. Explain why the parameter is not 0.8.
\end{enumerate}
\subsubsection*{Complete worked solution}
\par\noindent\begin{minipage}{\linewidth}
\textbf{Step 1: Define and identify units.} $X$ counts events over 2.5 hours; support $\{0,1,2,\ldots\}$ events. The count differs from rate and money.
\end{minipage}\par\medskip
\par\noindent\begin{minipage}{\linewidth}
\textbf{Step 2: Check and scale.} Under constant rate, independent disjoint-interval counts and isolated events, use Poisson. $\lambda=\text{rate}\times\text{exposure}=0.8\times2.5=2$ events.
\end{minipage}\par\medskip
\par\noindent\begin{minipage}{\linewidth}
\textbf{Step 3: Translate and calculate.} The event is $ X\geq3$. Symbolically,
\[P(X\geq a)=1-\sum_{j=0}^{a-1}e^{-\lambda}\lambda^j/j!.\]
Numerically,
\[P(X\geq3)=1-\sum_{j=0}^{2}e^{-2}(2)^j/j!\approx 0.3233.\]
\end{minipage}\par\medskip
\par\noindent\begin{minipage}{\linewidth}
\textbf{Step 4: Check and interpret.} The probability is in $[0,1]$; the full Poisson PMF sums to 1. The interval mean is 2 events; 0.8 is the hourly rate. The complement includes 0, 1 and 2, not 3.
\end{minipage}\par\medskip
\section{Notation reference and practice procedure}
\begin{center}
\begin{tabular}{lp{0.72\linewidth}}
\toprule
Symbol & Meaning and units \\ \midrule
$X,k$ & Event count and one non-negative integer value \\
$\lambda$ & Expected count in the selected exposure \\
rate & Expected events per unit exposure \\
$e$ & Mathematical constant, approximately 2.71828 \\
$k!$ & Factorial, with $0!=1$ \\
\bottomrule
\end{tabular}
\end{center}
\begin{enumerate}
\item Define the random variable in words and give its units and support.
\item List parameters and check that the model assumptions match the situation.
\item Translate ordinary language into an event, paying attention to endpoints.
\item Write the symbolic formula first, then substitute without rounding intermediate values.
\item Check bounds, total probability or a complementary calculation; interpret the answer in context.
\end{enumerate}
\Needspace{19\baselineskip}
\section{Foundational contextual practice}\label{sec:t2-c7-foundational}
\Needspace{14\baselineskip}
\subsection{No insurance claims}
Isolated claims occur at constant rate 0.6 per day, with independent counts in disjoint intervals.

\textbf{Questions/Tasks}
\begin{enumerate}
\item Calculate $P(X=0)$ for the exposure. 
\end{enumerate}
\subsubsection*{Complete worked solution}
\par\noindent\begin{minipage}{\linewidth}
\textbf{Step 1: Define and identify units.} $X$ counts events over 1 day; support $\{0,1,2,\ldots\}$ events. The count differs from rate and money.
\end{minipage}\par\medskip
\par\noindent\begin{minipage}{\linewidth}
\textbf{Step 2: Check and scale.} Under constant rate, independent disjoint-interval counts and isolated events, use Poisson. $\lambda=\text{rate}\times\text{exposure}=0.6\times1=0.6$ events.
\end{minipage}\par\medskip
\par\noindent\begin{minipage}{\linewidth}
\textbf{Step 3: Translate and calculate.} The event is $ X=0$. Symbolically,
\[P(X=k)=e^{-\lambda}\lambda^k/k!.\]
Numerically,
\[P(X=0)=e^{-0.6}(0.6)^0/0!\approx 0.5488.\]
\end{minipage}\par\medskip
\par\noindent\begin{minipage}{\linewidth}
\textbf{Step 4: Check and interpret.} The probability is in $[0,1]$; the full Poisson PMF sums to 1. For zero events, $\lambda^0=1$ and $0!=1$, so $P(X=0)=e^{-\lambda}$. 
\end{minipage}\par\medskip
\secInicio[][sec:t2-c7-foundational]
\Needspace{14\baselineskip}
\subsection{Two customer arrivals}
A desk has isolated arrivals at constant rate 3 customers per hour, with independent disjoint-interval counts.

\textbf{Questions/Tasks}
\begin{enumerate}
\item Calculate $P(X=2)$ for the exposure. 
\end{enumerate}
\subsubsection*{Complete worked solution}
\par\noindent\begin{minipage}{\linewidth}
\textbf{Step 1: Define and identify units.} $X$ counts events over 1 hour; support $\{0,1,2,\ldots\}$ events. The count differs from rate and money.
\end{minipage}\par\medskip
\par\noindent\begin{minipage}{\linewidth}
\textbf{Step 2: Check and scale.} Under constant rate, independent disjoint-interval counts and isolated events, use Poisson. $\lambda=\text{rate}\times\text{exposure}=3\times1=3$ events.
\end{minipage}\par\medskip
\par\noindent\begin{minipage}{\linewidth}
\textbf{Step 3: Translate and calculate.} The event is $ X=2$. Symbolically,
\[P(X=k)=e^{-\lambda}\lambda^k/k!.\]
Numerically,
\[P(X=2)=e^{-3}(3)^2/2!\approx 0.2240.\]
\end{minipage}\par\medskip
\par\noindent\begin{minipage}{\linewidth}
\textbf{Step 4: Check and interpret.} The probability is in $[0,1]$; the full Poisson PMF sums to 1. 
\end{minipage}\par\medskip
\secInicio[][sec:t2-c7-foundational]
\Needspace{19\baselineskip}
\section{Intermediate contextual practice}\label{sec:t2-c7-intermediate}
\Needspace{14\baselineskip}
\subsection{At most one defect}
Defects occur singly at constant rate 0.02 per metre, with independent counts on disjoint lengths. Inspect 50 metres.

\textbf{Questions/Tasks}
\begin{enumerate}
\item Calculate $P(X\leq1)$ for the exposure. 
\end{enumerate}
\subsubsection*{Complete worked solution}
\par\noindent\begin{minipage}{\linewidth}
\textbf{Step 1: Define and identify units.} $X$ counts events over 50 metres; support $\{0,1,2,\ldots\}$ events. The count differs from rate and money.
\end{minipage}\par\medskip
\par\noindent\begin{minipage}{\linewidth}
\textbf{Step 2: Check and scale.} Under constant rate, independent disjoint-interval counts and isolated events, use Poisson. $\lambda=\text{rate}\times\text{exposure}=0.02\times50=1$ events.
\end{minipage}\par\medskip
\par\noindent\begin{minipage}{\linewidth}
\textbf{Step 3: Translate and calculate.} The event is $ X\leq1$. Symbolically,
\[P(X\leq b)=\sum_{j=0}^{b}e^{-\lambda}\lambda^j/j!.\]
Numerically,
\[P(X\leq1)=\sum_{j=0}^{1} e^{-1}(1)^j/j!\approx 0.7358.\]
\end{minipage}\par\medskip
\par\noindent\begin{minipage}{\linewidth}
\textbf{Step 4: Check and interpret.} The probability is in $[0,1]$; the full Poisson PMF sums to 1. 
\end{minipage}\par\medskip
\secInicio[][sec:t2-c7-intermediate]
\Needspace{14\baselineskip}
\subsection{At least one payment alert}
Isolated alerts have constant rate 4 per hour and independent disjoint-interval counts. Observe 15 minutes.

\textbf{Questions/Tasks}
\begin{enumerate}
\item Calculate $P(X\geq1)$ for the exposure. Convert 15 minutes to hours.
\end{enumerate}
\subsubsection*{Complete worked solution}
\par\noindent\begin{minipage}{\linewidth}
\textbf{Step 1: Define and identify units.} $X$ counts events over 0.25 hours; support $\{0,1,2,\ldots\}$ events. The count differs from rate and money.
\end{minipage}\par\medskip
\par\noindent\begin{minipage}{\linewidth}
\textbf{Step 2: Check and scale.} Under constant rate, independent disjoint-interval counts and isolated events, use Poisson. $\lambda=\text{rate}\times\text{exposure}=4\times0.25=1$ events.
\end{minipage}\par\medskip
\par\noindent\begin{minipage}{\linewidth}
\textbf{Step 3: Translate and calculate.} The event is $ X\geq1$. Symbolically,
\[P(X\geq a)=1-\sum_{j=0}^{a-1}e^{-\lambda}\lambda^j/j!.\]
Numerically,
\[P(X\geq1)=1-\sum_{j=0}^{0}e^{-1}(1)^j/j!\approx 0.6321.\]
\end{minipage}\par\medskip
\par\noindent\begin{minipage}{\linewidth}
\textbf{Step 4: Check and interpret.} The probability is in $[0,1]$; the full Poisson PMF sums to 1. $15/60=0.25$ hours gives mean 1 alert. Multiplying an hourly rate by 15 minutes without conversion is a unit error.
\end{minipage}\par\medskip
\secInicio[][sec:t2-c7-intermediate]
\Needspace{19\baselineskip}
\section{Advanced contextual practice}\label{sec:t2-c7-advanced}
\Needspace{14\baselineskip}
\subsection{Claim costs over five days}
Isolated claims have constant rate 0.4 per day and independent disjoint-interval counts. Observe five days; each claim costs exactly 100 USD.

\textbf{Questions/Tasks}
\begin{enumerate}
\item Calculate $P(X\geq3)$ for the exposure. Interpret as a probability of claim cost at least 300 USD.
\end{enumerate}
\subsubsection*{Complete worked solution}
\par\noindent\begin{minipage}{\linewidth}
\textbf{Step 1: Define and identify units.} $X$ counts events over 5 days; support $\{0,1,2,\ldots\}$ events. The count differs from rate and money.
\end{minipage}\par\medskip
\par\noindent\begin{minipage}{\linewidth}
\textbf{Step 2: Check and scale.} Under constant rate, independent disjoint-interval counts and isolated events, use Poisson. $\lambda=\text{rate}\times\text{exposure}=0.4\times5=2$ events.
\end{minipage}\par\medskip
\par\noindent\begin{minipage}{\linewidth}
\textbf{Step 3: Translate and calculate.} The event is $ X\geq3$. Symbolically,
\[P(X\geq a)=1-\sum_{j=0}^{a-1}e^{-\lambda}\lambda^j/j!.\]
Numerically,
\[P(X\geq3)=1-\sum_{j=0}^{2}e^{-2}(2)^j/j!\approx 0.3233.\]
\end{minipage}\par\medskip
\par\noindent\begin{minipage}{\linewidth}
\textbf{Step 4: Check and interpret.} The probability is in $[0,1]$; the full Poisson PMF sums to 1. $L=100X$ USD gives $L\geq300$ exactly when $X\geq3$. This translation relies on the fixed cost per claim.
\end{minipage}\par\medskip
\secInicio[][sec:t2-c7-advanced]
\Needspace{14\baselineskip}
\subsection{Defaults and a clustering warning}
Version A has isolated defaults at constant rate 0.5 per month over stable exposure, with independent disjoint-interval counts. Observe two months. In version B, a market shock triggers many defaults together.

\textbf{Questions/Tasks}
\begin{enumerate}
\item Calculate $P(X\leq1)$ for the exposure. Calculate for A; explain why rarity alone does not justify the model for B.
\end{enumerate}
\subsubsection*{Complete worked solution}
\par\noindent\begin{minipage}{\linewidth}
\textbf{Step 1: Define and identify units.} $X$ counts events over 2 months; support $\{0,1,2,\ldots\}$ events. The count differs from rate and money.
\end{minipage}\par\medskip
\par\noindent\begin{minipage}{\linewidth}
\textbf{Step 2: Check and scale.} Under constant rate, independent disjoint-interval counts and isolated events, use Poisson. $\lambda=\text{rate}\times\text{exposure}=0.5\times2=1$ events.
\end{minipage}\par\medskip
\par\noindent\begin{minipage}{\linewidth}
\textbf{Step 3: Translate and calculate.} The event is $ X\leq1$. Symbolically,
\[P(X\leq b)=\sum_{j=0}^{b}e^{-\lambda}\lambda^j/j!.\]
Numerically,
\[P(X\leq1)=\sum_{j=0}^{1} e^{-1}(1)^j/j!\approx 0.7358.\]
\end{minipage}\par\medskip
\par\noindent\begin{minipage}{\linewidth}
\textbf{Step 4: Check and interpret.} The probability is in $[0,1]$; the full Poisson PMF sums to 1. Version B may have clustered events and dependent counts. Low frequency does not repair these failures; the numerical answer applies to A.
\end{minipage}\par\medskip
\secInicio[][sec:t2-c7-advanced]
\section{Closing reference: key ideas and common errors}
Fix the interval and units first. Multiply rate by exposure for $\lambda$. Include zero and use $0!=1$. At least $a$ complements counts through $a-1$. Check stable rate, independent intervals and isolated events; rarity alone is insufficient. Money is a separate variable.
\secInicio
\end{document}
